\chapter{Происхождение спектральной плотности резервуара для логарифма определителя KMS}
\label{ch:v9-kms-reservoir-spectral-density-origin}

\section{Вопрос происхождения}

Keldysh admission показал, что conductance может входить в dissipative
kernel, но normalized contour не оставляет требуемый vacuum логарифм определителя. Теперь
проверяется более сильная гипотеза: определяют ли три on-shell conductances
полный reservoir spectral profile, а вместе с ним off-shell self-energy и
bath determinant.

Для трёх type-gaps возьмём exact witness
\begin{equation}
 \Delta_s=1,
 \qquad \Delta_a=2,
 \qquad \Delta_t=3,
 \qquad m=(1,1,3).
 \label{eq:v9-kms-reservoir-gaps}
\end{equation}
Conductances считывают только on-shell values:
\begin{equation}
 \kappa_\alpha=J(\Delta_\alpha),
 \qquad \alpha\in\{s,a,t\}.
 \label{eq:v9-kms-reservoir-onshell-rates}
\end{equation}

\section{Evaluation map и его kernel}

На полиномах степени не выше шести evaluation map имеет matrix
\begin{equation}
 E=\begin{pmatrix}
 1&1&1&1&1&1&1\\
 1&2&4&8&16&32&64\\
 1&3&9&27&81&243&729
 \end{pmatrix}.
 \label{eq:v9-kms-reservoir-evaluation-map}
\end{equation}
Его rank/nullity равны
\begin{equation}
 \rank E=3,
 \qquad \dim\ker E=7-3=4.
 \label{eq:v9-kms-reservoir-evaluation-rank}
\end{equation}
Следовательно, даже в finite polynomial class три rates оставляют четыре
off-shell направления.

Явный положительный kernel-profile строится из
\begin{equation}
 q(\omega)=
 (\omega-1)^2(\omega-2)^2(\omega-3)^2\geq0.
 \label{eq:v9-kms-reservoir-kernel-profile}
\end{equation}
Он точно невидим на gaps:
\begin{equation}
 q(1)=q(2)=q(3)=0.
 \label{eq:v9-kms-reservoir-kernel-zeros}
\end{equation}

\section{Два положительных reservoir profiles}

На interval $0\leq\omega\leq4$ положим
\begin{equation}
 J_0(\omega)=1,
 \qquad
 J_1(\omega)=1+\frac1{16}q(\omega).
 \label{eq:v9-kms-reservoir-two-profiles}
\end{equation}
Оба profile строго положительны:
\begin{equation}
 J_0(\omega)=1>0,
 \qquad J_1(\omega)\geq1>0.
 \label{eq:v9-kms-reservoir-positivity}
\end{equation}
Их on-shell conductance vectors совпадают:
\begin{equation}
 \boldsymbol\kappa[J_0]
 =\boldsymbol\kappa[J_1]
 =(1,1,1)^T.
 \label{eq:v9-kms-reservoir-equal-rates}
\end{equation}
Weighted type normalization также совпадает:
\begin{equation}
 \kappa_s+\kappa_a+3\kappa_t=5.
 \label{eq:v9-kms-reservoir-equal-normalization}
\end{equation}

\section{Off-shell moment defect}

Zeroth spectral moments различны:
\begin{equation}
 \mu_0[J_0]=\int_0^4J_0(\omega)d\omega=4,
 \qquad
 \mu_0[J_1]=\frac{527}{105}.
 \label{eq:v9-kms-reservoir-zeroth-moments}
\end{equation}
Точный defect равен
\begin{equation}
 \mu_0[J_1]-\mu_0[J_0]=\frac{107}{105}>0.
 \label{eq:v9-kms-reservoir-zeroth-defect}
\end{equation}
Первые moments также различны:
\begin{equation}
 \mu_1[J_0]=8,
 \qquad
 \mu_1[J_1]=\frac{1054}{105},
 \label{eq:v9-kms-reservoir-first-moments}
\end{equation}
причём
\begin{equation}
 \mu_1[J_1]-\mu_1[J_0]=\frac{214}{105}>0.
 \label{eq:v9-kms-reservoir-first-defect}
\end{equation}

Для large complex $z$ reservoir self-energy имеет expansion
\begin{equation}
 \Sigma_J(z)=\int_0^4\frac{J(\omega)}{z-\omega}\,d\omega
 =\frac{\mu_0[J]}{z}+\frac{\mu_1[J]}{z^2}+O(z^{-3}).
 \label{eq:v9-kms-reservoir-self-energy-expansion}
\end{equation}
Поэтому equal on-shell rates не определяют даже первые два off-shell
коэффициента:
\begin{equation}
 \Sigma_{J_1}(z)-\Sigma_{J_0}(z)
 =\frac{107}{105z}+\frac{214}{105z^2}+O(z^{-3})\neq0.
 \label{eq:v9-kms-reservoir-self-energy-defect}
\end{equation}
Bath логарифм определителя зависит от полного analytic self-energy, поэтому не является
функцией только трёх значений $\kappa_\alpha$.

\section{Type strengths и normalization}

Если допустить independent type strengths
\begin{equation}
 J_\alpha(\omega)=c_\alpha J(\omega),
 \qquad c_\alpha>0,
 \label{eq:v9-kms-reservoir-type-strengths}
\end{equation}
то одна weighted normalization имеет rank one:
\begin{equation}
 A_{\rm norm}=\begin{pmatrix}1&1&3\end{pmatrix},
 \qquad \rank A_{\rm norm}=1,
 \qquad \dim\ker A_{\rm norm}=2.
 \label{eq:v9-kms-reservoir-normalization-rank}
\end{equation}
Значит, reservoir ansatz не выбирает ещё и два relative type strengths.

\section{ProofDSL и итог}

LCF certificate проверяет evaluation rank, exact kernel vector, equal
rates, moments, positive defects и off-shell nonconstancy:
\begin{equation}
 N_{\rm ProofDSL\ obligations}=12/12,
 \qquad N_{\rm registered\ gates}=59,
 \qquad N_{\rm registered\ obligations}=480.
 \label{eq:v9-kms-reservoir-proofdsl}
\end{equation}
Итоговые ledgers:
\begin{equation}
 \begin{aligned}
 N_{\rm positive\ spectral\ profiles}&=2/2,\\
 N_{\rm equal\ on\ shell\ rates}&=3/3,\\
 N_{\rm unique\ off\ shell\ profile}&=0/1,\\
 N_{\rm rates\ determine\ logdet}&=0/1,\\
 N_{\rm reservoir\ parent\ origin}&=0/1,\\
 N_{\rm physical\ four\ slot\ parent}&=0/1.
 \end{aligned}
 \label{eq:v9-kms-reservoir-ledger}
\end{equation}

\begin{theorem}[On-shell reproduction]
Positive reservoir profiles могут точно воспроизвести все три KMS
conductances и их weighted normalization.
\end{theorem}

\begin{nogocriterion}[Off-shell spectral nonuniqueness]
Три on-shell values имеют четырёхмерное kernel уже на полиномах степени
шесть. Два strictly positive exact profiles дают одинаковые rates, но
разные zeroth/first moments и разные self-energy asymptotics. Поэтому
conductances не определяют спектральная плотность резервуара или bath логарифм определителя.
\end{nogocriterion}

\begin{protocol}\begin{enumerate}
\item polynomial profile dimension: \textbf{$7$};
\item on-shell map rank/nullity: \textbf{$3/4$};
\item rate vectors: \textbf{$(1,1,1)$ и $(1,1,1)$};
\item zeroth moment defect: \textbf{$107/105$};
\item first moment defect: \textbf{$214/105$};
\item relative type-strength freedom: \textbf{$2$};
\item reservoir/логарифм определителя origin: \textbf{$0/1$};
\item ProofDSL: \textbf{$12/12$, registry $59/480$};
\item следующий гейт: \textbf{reservoir measure-anomaly parent-origin}.
\end{enumerate}\end{protocol}

Машинный сертификат сохранён в
\path{s2t_v9_endpoint_creation_kms_logdet_reservoir_spectral_density_parent_origin_gate_results.json}.

\begin{thebibliography}{9}
\bibitem{v9-kms-reservoir-davies}
E.~B.~Davies,
\emph{Markovian master equations},
Commun. Math. Phys. \textbf{39} (1974) 91--110.
\bibitem{v9-kms-reservoir-tedopa}
A.~N"usseler, I.~Dhand, S.~F.~Huelga, M.~B.~Plenio,
\emph{Efficient simulation of open quantum systems coupled to a fermionic bath},
Phys. Rev. B \textbf{101} (2020) 155134, arXiv:1909.09589.
\end{thebibliography}